%==========================================================================================
\section{Notebooks}\label{sec:notebooks}
%==========================================================================================

The notebooks chain the modules into complete analyses. They are not documented function by
function: each is organised in numbered sections with explanations, and its settings are
gathered in its first cells. Table~\ref{tab:notebooks} summarises what each one does and which
functions it relies on.

\begin{table}[h]
\centering\small
\caption{Notebooks of the repository.}
\label{tab:notebooks}
\begin{tabularx}{\textwidth}{@{}l X@{}}
\toprule
Notebook & Content and main functions \\
\midrule
\fn{fluo_analysis_clean.ipynb} & The white-beam workflow on a BLISS mesh scan:
\fn{BlissScan}; normalisation (\fn{normalise_spectra}); surface region (the whole map for a
film, a facet of a crystal: \fn{SurfaceRegion}, \fn{RegionSelector}); map-averaged spectrum and
counting error; peak scouting (\fn{identify_peaks}); BM32 beam from the scan (\fn{bm32_from_h5}),
detector and sample matrix (\fn{MatrixModel}); internal energy calibration starting from the
C10 calibration; pile-up and escape peaks; line groups and constrained fit
(\fn{element_line_groups}, \fn{fit_average_spectrum}); sensitivity to the element list;
resolution stability (\fn{width_window}, \fn{peak_width_profile}, \fn{detect_steps},
\fn{fit_segments}); per-pixel maps (\fn{fit_map}) with quality report, individual pixel
fits, quantitative ratios (\fn{quantify_layer}, \fn{atom_ratio}) and closure tests. \\
\fn{fluo_analysis.ipynb} & Earlier version of the same workflow, kept with its outputs. It
does not include the internal calibration step. \\
\fn{xrf_energy_calibration.ipynb} & Energy calibration on the AXO RF C10 reference
(\fn{calibrate}, \fn{plot_calibration}), with the beamline settings read from the scan, the
fitted air path, line checks and intensity check, comparison with the calibration in use, and
a JSON record that the analysis notebooks read as their starting calibration. \\
\fn{xrf_mono_spectrum.ipynb} & A spectrum excited by a monochromatised Mo tube (the PyMca
training spectrum \fn{XRFSpectrum.mca}, not distributed with the repository; see the PyMca
tutorial): calibration from two peaks (\fn{calibration_from_peaks}), peak scouting, first-pass
residuals revealing Tl and Pb, full fit up to the elastic peak with detector tails and the
impulse-approximation Compton peak (\fn{ScatterSource}, \fn{fit_tails}), rejected candidates,
and concentrations from the Co internal standard. \\
\bottomrule
\end{tabularx}
\end{table}

\paragraph{Results files.} The notebooks write their results to \fn{processed/}: fit records
(JSON), average spectra and maps (NPZ), and figures. The fit records hold the calibration, the
detector parameters, every fitted parameter with its error and the group areas, so that later
notebooks can start from them.
